\begin{abstract}
We build one model that classifies a synthetic aperture radar (SAR) image into one of nine scene classes and detects
every object in it, within 130\,M parameters (40\,M trainable).
Three properties of the data shaped the design.
Half of the boxes are smaller than a 16-pixel ViT patch, so instead of adapting the reference ViT-B/16 with LoRA we
fully fine-tune a COCO-pretrained YOLO26-L detector (26.2\,M parameters, stride-8 features) and attach a 5.6\,k-parameter
classifier to its trunk.
Two thirds of held-out training images have a near-duplicate in the rest of the training set, so we validate on a
holdout that mirrors this and favour augmentations that keep scene geometry: SAR speckle and gamma noise gain +2.3 to
+2.9 $\Delta_m$ over three seeds, while zoom and mosaic lose about 5.
All boxes of an image share its class, which a parameter-free cross-task decoding exploits: detector evidence refines
the image class and the class re-ranks the boxes (+0.61\,$\pm$\,0.15 $\Delta_m$ over nine runs).
With weight averaging, the final single model reaches $\Delta_m$ +45.9 on val and +48.7 on test, against +4.4 and +3.8
for the best in-budget baseline, at 2.5\,ms per image.
\end{abstract}
